\documentclass[a4paper]{article}
\usepackage{graphicx}
\usepackage{amsmath,amssymb}
\usepackage{minted}
\usepackage{multirow}
\usepackage{float}
\usepackage{algorithm}
\usepackage{algpseudocode}
\usepackage{todonotes}
\usepackage{forest}
\usepackage{xstring}
\usepackage{xcolor}
\usepackage[most]{tcolorbox}
\usepackage{xparse}
\usepackage{natbib}

\title{Large Source document virtualization}
\date{2026-09-03}
\author{Mohammad Rahmani}

\begin{document}
\maketitle

\section{Bug}
Opening a multi-megabyte LaTeX document could freeze the Qt GUI and eventually
terminate Mind Mirror with exit code $-9$. Visual projection virtualization and
asynchronous Structure parsing were insufficient because the complete Source
was still materialized in a \texttt{QTextEdit}/\texttt{QTextDocument}.

\section{Fix}
A canonical \texttt{LatexSourceBuffer} now owns the complete LaTeX source and
exposes a bounded, line-aligned Source window. \texttt{LargeLatexEditor}
subclasses the normal \texttt{LatexEditor}; small documents retain the normal
behavior, while large documents place only the bounded Source window in Qt.
Edits are spliced exactly once into the canonical source and the line index is
updated incrementally.

Save, autosave, Preview, Visual, and Structure continue receiving the complete
canonical source. Structure navigation maps full-document line numbers back to
the correct Source window, so the complete Structure remains available even for
large documents.

\section{Regression protection}
Tests verify that the canonical source remains complete, an edited Source
window is spliced only once without duplicating the untouched suffix, global
line mapping survives window edits, and the editor-panel path uses the OOP
large-document editor specialization without runtime method patching.

\end{document}
